\section{Formalization}
\label{sec:blockchain}

An Automated Market Maker implements a decentralized exchange between two different token types. The exchange rate is determined by a smart contract, which also takes care of performing the exchange itself: namely, the contract receives from a trader some amount of the input token type, and sends back the correct amount of the output token type, which is taken from the AMM reserves. A single smart contract can control many instances of AMMs (also called \emph{AMM pairs}): we may have a pair for each possible unordered pair of token types. To create an AMM instance, a user must provide the initial liquidity for the reserves of that pair of tokens. Liquidity providers are rewarded with a type of token that specifically represents shares in that AMM's reserves: we call these \emph{minted} token types, while any other token type will be called \emph{atomic}.

\paragraph*{Blockchain State}

We begin by formalizing the blockchain state, abstracting from the details that are immaterial to the study of AMMs. 
Then, our model includes the users' wallets, the AMMs and their reserves
(see~\Cref{tab:amm-system}).
We formalize the universes of users and atomic token types as the types \code{A} and \code{T}, respectively, as structures that encapsulate a natural number.
Hereafter, we use $a, b, \ldots$ to denote users in $\code{A}$, 
and $\tau, \tau_0, \tau_1, \ldots$ to denote tokens in $\code{T}$.
Minted token types are pairs of \code T. 
We represent the funds owned by a user by a \emph{wallet} that maps token types to non-negative reals.
%
To rule out wallets with infinite tokens,
we use Mathlib's finitely supported functions\footnote{\url{https://leanprover-community.github.io/mathlib4_docs/Mathlib/Data/Finsupp/Defs.html}}: 
in general,
given any type $\alpha$ and any type $M$ with a $0$ element, $f \in \alpha \to_0 M$ if $\mathrm{supp}\left(f\right) = \left\{x \in \alpha \mid f\left(x\right) \neq 0\right\}$ is finite.  

We define the type $\code W_0$ of wallets of atomic tokens as a structure encapsulating $\code T \to_0 \mathbb R_{\geq 0}$. This definition induces an element $0 \in \code W_0$ such that $\mathrm{supp}\left(0\right) = \emptyset$: this is the \emph{empty wallet}, which enables us to form the type $\code T \to_0 \code W_0$. 
%
We define the type $\code W_1$ of wallets of minted tokens 
as a structure that encapsulates a function $\code{bal} \in \code T \to_0 \code{W_0}$. 
The intuition is that $\code{bal}\ \tau_0\ \tau_1$ gives the owned amount of the minted token type created by the AMM pair with tokens $\tau_0$ and $\tau_1$.
Consistently, the function \code{bal} must satisfy two conditions: 
$\code{bal}\ \tau_0\ \tau_1 = \code{bal}\ \tau_1\ \tau_0$, 
meaning that the order of atomic tokens is irrelevant, 
and $\code{bal}\ \tau\ \tau = 0$, meaning that the two token types 
in an AMM must be distinct. 
Our definition of $\code W_1$ encapsulates proofs of these two properties, called \code{unord} and \code{distinct}, respectively.

We map users to their wallets with the types $\code S_0$ and $\code S_1$,
which account for the atomic tokens and for the minted tokens,
respectively.
Finally, we formalize sets of AMM pairs with the type \code{AMMs}. 
The definition is strikingly similar to that of $\code W_1$, but with a changed constraint. 
The intuition is that $\code{res}\ \tau_0\ \tau_1$ gives the reserves of $\tau_0$ in the AMM pair $(\tau_0$, $\tau_1)$,
while  $\code{res}\ \tau_1\ \tau_0$ gives the reserves of $\tau_1$ in the same AMM pair. 
For uninitialized AMM pairs, the obtained reserves must be $0$. 
The property \code{posres} ensures that either an AMM pair has 
no reserves of both token types (\ie, the AMM has not been created yet),
or both token types have strictly positive reserves
(\ie, one cannot deplete the reserves of a single token type in an AMM pair).
% 
We combine the previous definitions in the type $\Gamma$, which represents the state of a blockchain (note that $\Gamma$ abstracts from all the details immaterial for AMMs). 

\paragraph*{Token supply}

Given a blockchain state $s \in \Gamma$, we define the supply of an atomic token type $\tau_0$ as:
\[
\code{atomsupply}_s\left(\tau_0\right) 
\;\; = \;\; 
\sum_{a \in \mathrm{supp}(s\code{.atoms})} 
\hspace{-15pt}
(s\code{.atoms} \; a) \; \tau_0 
\;\; + \;\;
\sum_{\tau_1 \in \mathrm{supp}(s\code{.amms} \, \tau_0)} 
\hspace{-20pt}
(s\code{.amms} \; \tau_0) \; \tau_1
\]
where the partial application $s\code{.amms} \; \tau_0$ 
gives a map with all AMM pairs with $\tau_0$ as one of their token types.
We define the supply of a minted token type $(\tau_0,\tau_1)$ as follows:
\[
\code{mintsupply}_s\left(\tau_0, \tau_1\right) 
\;\; = \;\; 
\sum_{a \in \mathrm{supp}(s\code{.mints})} 
\hspace{-15pt}
(s\code{.mints} \; a) \; \tau_0 \;\tau_1
\]

The corresponding Lean definitions (in \Cref{list:supplydefs}) 
have been split in order to facilitate theorem proving and, in particular, the use of Lean's simplifier.

\begin{table}[t]
    \caption{Fundamental Lean definitions for the state of an AMM system.}
    \label{tab:amm-system}
    \vspace{-20pt}
    \centering
    \begin{tabular}{p{0.25\linewidth}p{0.30\linewidth}p{0.45\linewidth}}
    \begin{lstlisting}[numbers=none]
structure A where
  n: ℕ

structure T where
  n: ℕ

structure W₀ where
  bal: T →₀ ℝ≥0
  
structure S₀ where
  map: A →₀ W₀\end{lstlisting} & 
    \begin{lstlisting}[numbers=none]
structure W₁ where
  bal: T →₀ W₀
  unord: ∀ (τ₀ τ₁: T), 
    bal τ₀ τ₁ = f τ₁ τ₀
  distinct: ∀ (τ: T), 
    bal τ τ = 0
  
structure S₁ where
  map: A →₀ W₁ \end{lstlisting} &
    \begin{lstlisting}[numbers=none]
structure AMMs where
  res: T →₀ W₀
  distinct: ∀ (τ: T), 
    res τ τ = 0
  posres: ∀ (τ₀ τ₁: T), 
    res τ₀ τ₁ ≠ 0  ↔ f τ₁ τ₀ ≠ 0

structure Γ where
  atoms: S₀
  mints: S₁
  amms: AMMs\end{lstlisting}
    \end{tabular}
\end{table}

\begin{table}[t]
\caption{Supply of atomic token types and of minted token types.}
\label{list:supplydefs}
\begin{lstlisting}
noncomputable def S₀.supply (s: S₀) (τ: T): ℝ≥0 := s.map.sum (λ _ w => w τ)

noncomputable def S₁.supply (s: S₁) (τ₀ τ₁: T): ℝ≥0 := 
    s.map.sum (λ _ w => w.get τ₀ τ₁)

noncomputable def AMMs.supply (amms: AMMs) (τ: T): ℝ≥0 := 
    (amms.res τ).sum λ _ x => x

noncomputable def Γ.atomsupply (s: Γ) (τ: T): ℝ≥0 := 
    (s.atoms.supply τ) + (s.amms.supply τ)

noncomputable def Γ.mintsupply (s: Γ) (τ₀ τ₁: T): ℝ≥0 := s.mints.supply τ₀ τ₁
\end{lstlisting}
\end{table}
\paragraph*{AMM reserves}

Given a blockchain state $s$ and two token types $\tau_0$, $\tau_1$,
the terms $s.\code{amms} \; \tau_0 \; \tau_1$
and $s.\code{amms} \; \tau_1 \; \tau_0$
denote, respectively, the reserves of $\tau_0$ and $\tau_1$
in the AMM pair $(\tau_0,\tau_1)$.
%
This way of accessing the AMM reserves is a bit impractical:
when writing proofs, using $s.\code{amms} \; \tau_0 \; \tau_1$
carries an obligation to provide the functions $\code{distinct}$ and $\code{posres}$.
In particular, this requires the user to explicitly add, in any theorem
using the reserves, the assumption that the reserves
are strictly positive to indicate the AMM pair has been created.
Furthermore, this way of accessing the reserves hides the
fact that when one of the reserves is strictly positive, 
also the other one is such, which again should be made explicit when
writing proofs.

To cope with these issues, 
we build a Lean API that allows for hiding these implementation details (see~\Cref{list:ammapi}). 
For example, given an AMM pair $(\tau_0, \tau_1)$ in a state $s$, 
the expression $s\code{.amms.r_0 \, \tau_0 \; \tau_1 \; init}$ 
give the reserves of token $\tau_0$ in the AMM, 
which are guaranteed to be strictly positive 
under the initialization precondition $\code{init} \in (s\code{.amms.init})$.

\begin{table}[t]
\caption{Fragment of the AMM API: AMM reserves.}
\label{list:ammapi}
\begin{lstlisting}
def AMMs.init (amms: AMMs) (τ₀ τ₁: T): Prop := amms.res τ₀ τ₁ ≠ 0

def AMMs.r₀ (amms: AMMs) (τ₀ τ₁: T) (h: amms.init τ₀ τ₁): ℝ>0 := ⟨ amms.res τ₀ τ₁, 
    by unfold init at h; exact NNReal.neq_zero_imp_gt h ⟩

def AMMs.r₁ (amms: AMMs) (τ₀ τ₁: T) (h: amms.init τ₀ τ₁): ℝ>0 := ⟨ amms.res τ₁ τ₀, 
    by unfold init at h; exact NNReal.neq_zero_imp_gt ((amms.posres τ₀ τ₁).mp h) ⟩
\end{lstlisting}
\end{table}

\paragraph*{Transactions}

Our model encompasses all the main types of transactions supported by AMMs: creating an AMM, adding/removing liquidity, and swapping a token for another. 
Swaps are parameterised by a \emph{swap rate function}, which determines the  exchange rate.
We use the formalization of swap transactions (\Cref{list:swapdefs}) to exemplify the scheme we used for all the transaction types. 

The type $\code{Swap}\left(sx, s, a, \tau_0, \tau_1, x\right)$ represents valid swap transactions in a blockchain state $s$, with the swap rate function $sx$, performed by user $a$ to exchange $x$ amount of the input token $\tau_0$ for a certain amount of the output token $\tau_1$ (\Cref{lst:swapsig}). 
Each element of this type is a structure containing a proof of the validity of the transaction. 
For example, for swap transactions we must prove that the user has enough amount of $\tau_0$ (condition \code{enough}), 
that the AMM pair with tokens $\tau_0$ and $\tau_1$ exists 
(condition \code{exi}), 
and it has enough reserves of $\tau_1$ to give as output 
(condition \code{nodrain}). 
%
Since the type $\code{Swap}\left(\cdots\right)$
is empty when the parameters do not satisfy the above conditions,
invalid transactions are not really expressible in our model. 
%
Instead, if $\code{Swap}\left(\cdots\right)$ represents a valid transaction, it will be a singleton type due to proof irrelevance 
(\ie, any two proofs of the same proposition are equal). 

Each transaction is equipped with an $\code{apply}$ function that yields the state reached by executing the transaction in the given state. 
For example, for $sw \in \code{Swap}\left(sx, s, a, \tau_0, \tau_1, x\right)$, 
$\code{apply}\; sw$ yields a state where:
\begin{inlinelist}
\item $a$'s atomic tokens wallet is updated by removing $x$ units of $\tau_0$ and adding $sw\code{.y}$ units of $\tau_1$ (\Cref{lst:updatoms}),
where $sw\code{.y}$ is the amount of tokens outputted by the AMM pair; 
\item accordingly, the AMM reserves are updated by removing $sw\code{.y}$ units of $\tau_1$ and adding $x$ units of $\tau_0$ (\Cref{lst:suby}); 
\item the minted token wallets is unchanged (\Cref{lst:unchangedmints}).
\end{inlinelist}
These definitions use functions and proofs 
not included in \Cref{list:swapdefs} for brevity, 
such as \code{sub} and \code{sub\_r1}. 
These are designed with the same spirit of allowing only valid operations, and so require suitable proofs.
For example, \code{sub\_r1} requires a proof of the existence of the AMM we are removing liquidity from, and a proof that the AMM pair has enough liquidity to retain a positive amount of reserves. 
We build these proofs inline using those contained in the structure:
for instance, the parameter $sw\code{.exi}$ passed to \code{sub\_r1} 
at line 11 is a proof that the AMM pair exists. 
Then, at line 16 we define the constant-product swap rate function, that is the swap rate function used by Uniswap v2. From~\Cref{th:swapdirection} onwards, our results will focus on AMMs using this swap rate function.

\begin{table}[t]
\caption{Definition of the swap transaction type and of its application, as well as the constant-product swap rate function.}
\label{list:swapdefs}
\begin{lstlisting}
abbrev SX := ℝ>0 → ℝ>0 → ℝ>0 → ℝ>0
structure Swap (sx: SX) (s: Γ) (a: A) (τ₀ τ₁: T) (x: ℝ>0) where @\label{lst:swapsig}@
  enough: x ≤ s.atoms.get a τ₀      -- user a has at least x τ₀
  exi: s.amms.init τ₀ τ₁            -- AMM pair τ₀ τ₁ exists in s
  nodrain: x*(sx x (s.amms.r₀ τ₀ τ₁ exi) (s.amms.r₁ τ₀ τ₁ exi))
           < (s.amms.r₁ τ₀ τ₁ exi)  -- AMM has enough output tokens

def Swap.y (sw: Swap sx s a τ₀ τ₁ x): ℝ>0 := 
    x*(sx x (s.amms.r₀ τ₀ τ₁ sw.exi) (s.amms.r₁ τ₀ τ₁ sw.exi))

noncomputable def Swap.apply (sw: Swap sx s a τ₀ τ₁ x): Γ := {
  atoms := (s.atoms.sub a τ₀ x sw.enough).add a τ₁ sw.y, @\label{lst:updatoms}@
  mints := s.mints, @\label{lst:unchangedmints}@
  amms  := (s.amms.sub_r₁ τ₀ τ₁ sw.exi sw.y sw.nodrain).add_r₀ τ₀ τ₁ (by simp[sw.exi]) x } @\label{lst:suby}@

noncomputable def SX.constprod: SX := λ (x r₀ r₁: ℝ+) => r₁/(r₀ + x)
\end{lstlisting}
\end{table}

\paragraph*{Price, networth and Gain}

An important aim of our model is to state and prove economic properties of AMMs related to the networth of their users. The fundamental definitions are in~\Cref{list:networth}. Given a wallet $w \in \code W_0$ and an atomic token price oracle $o \in \code T \to \mathbb R_{>0}$, we define the \emph{value} of $w$ in \cref{lst:atomicworthdef} as:
\[
\code{value}\left(w, o\right) = \sum\limits_{\tau \in \mathrm{supp}\left(w\right)} w\left(\tau\right) \cdot o\left(\tau\right)
\] 
The value of a wallet of minted tokens $w \in \code W_1$
is defined similarly, except that: 
\begin{inlinelist}
\item the summation ranges over $\mathrm{supp}\left(w\code{.u}\right)$, with $w\code{.u} \in \code T^2 \to_0 \mathbb R_{\geq 0}$ representing the uncurrying of $w$;
\item the summation is divided by 2 since, if $\left(\tau_0, \tau_1\right)$ is in the support of $w$, then also $\left(\tau_1, \tau_0\right)$ is in its support.
\end{inlinelist}
For uniformity with the definition of value of atomic wallets, 
also here we assume an oracle that gives the price of (minted) tokens. 
However, while for pricing atomic tokens we indeed resort to an oracle,
for minted tokens this oracle is instantiated to a specific function,
coherently with~\cite{Bartoletti_2022}: 
\[
\code{mintedprice}_s(o, \tau_0, \tau_1) = \frac
    {(s.\code{amms.r_0} \; \tau_0 \; \tau_1) \cdot o(\tau_0) + (s.\code{amms.r_1} \; \tau_0 \; \tau_1) \cdot o(\tau_1)}
    {\code{mintsupply}_s(\tau_0,\tau_1)}
\]
where we have omitted the initialization precondition $\code{h}$ for brevity.

We then define the \emph{networth} of a user as the sum of the value of their two types of wallets (\cref{lst:networthsig}). 
The \emph{gain} of a user upon an update of the blockchain state is the difference between the networth in the new state and that in the old state (\cref{lst:gaindef}).

\begin{table}[t]
\caption{Users' networth and gain.}
\label{list:networth}
% [caption={Fundamental definitions for a user's networth.},label={list:networth},captionpos=t,float,abovecaptionskip=-\medskipamount]
\begin{lstlisting}
noncomputable def W₀.value (w: W₀) (o: T → ℝ>0): ℝ≥0 := @\label{lst:atomicworthsig}@
  w.sum (λ τ x => x*(o τ)) @\label{lst:atomicworthdef}@

noncomputable def W₁.value (w: W₁) (o: T → T → ℝ≥0): ℝ≥0 :=
  (w.u.sum (λ p x => x*(o p.fst p.snd))) / 2

noncomputable def Γ.mintedprice (s: Γ) (o: T → ℝ>0) (τ₀ τ₁: T): ℝ≥0 := @\label{lst:mintedpricesig}@
  if h:s.amms.init τ₀ τ₁ then
  ((s.amms.r₀ τ₀ τ₁ h)*(o τ₀)+(s.amms.r₁ τ₀ τ1 h)*(o τ₁)) / (s.mints.supply τ₀ τ₁)
  else 0 -- price is zero if AMM is not initialized

noncomputable def Γ.networth (s: Γ) (a: A) (o: T → ℝ>0): ℝ≥0 := @\label{lst:networthsig}@
  (W₀.value (s.atoms.get a) o) + (W₁.value (s.mints.get a) (s.mintedprice o))

noncomputable def A.gain (a: A) (o: T → ℝ>0) (s s': Γ): ℝ :=
  ((s'.networth a o): ℝ) - ((s.networth a o): ℝ) @\label{lst:gaindef}@
\end{lstlisting}
\end{table}


\paragraph*{Reachable states}

To formalize reachable states, we begin by defining sequences of transactions (\cref{tab:reachable}). 
$\code{Tx}\left(sx,s,s'\right)$ is the type of sequences of valid transactions
(of any kind)
starting from state $s$ and leading to $s'$.
The parameter $sx$ is the swap rate function used in swap transactions.
Technically, $\code{Tx}\left(sx,s,s'\right)$
is an instance of the indexed family of dependent types \code{Tx}, dependent on $sx$ and $s$, and indexed by $s'$. In practice, this means that the constructors must preserve the values of $sx$ and $s$ (building upon the sequence of transactions does not change the swap rate function being used nor the originating state), while $s'$ may change after each construction (the state resulting from the sequence changes with each transaction that is added to it). 
A state $s'$ is \emph{reachable} if there exists a valid sequence of transactions that reaches $s'$ starting from a valid \emph{initial} state $s$, \ie a state with no initialized AMMs or minted token types in circulation.

\begin{table}[t]
    \caption{Sequences of transactions and reachable states.}
    % \vspace{-10pt}
    \label{tab:reachable}
    \begin{lstlisting}
inductive Tx (sx: SX) (init: Γ): Γ → Type where @\label{lst:txdef}@
  | empty: Tx sx init init

  | swap (s': Γ) (rs: Tx sx init s')
         (sw: Swap sx s' a τ₀ τ₁ v0):
      Tx sx init sw.apply
  -- Other constructors omitted for brevity
  
def validInit (s: Γ): Prop :=
  (s.amms = AMMs.empty ∧ s.mints = S₁.empty)

def reachable (sx: SX) (s: Γ): Prop :=
  ∃ (init: Γ) (tx: Tx sx init s), validInit init
\end{lstlisting}
\end{table}